\ifdefined\gkver\else\def\gkver{student}\fi
\documentclass[\gkver]{gaokaozhenti}
\newcommand{\ccnum}[1]{{\fontspec{Noto Sans CJK JP}#1}}
\begin{document}

\chapter{2005年全国理综I}
%% paperType: 理综物理部分

\begin{enumerate}
\item
%% number: 14
%% paperName: 2005年全国理综I第 14 题 6 分
%% typeId: 1
%% type: 单选题
%% chapter: 运动和力
%% point: 超重失重
%% method: 无
%% score: 6
%% degree: 750
%% duplicateId: 0
%% body:
一质量为 $m$ 的人站在电梯中，电梯加速上升，加速度大小为 $\frac{1}{3}g$，$g$ 为重力加速度。人对电梯底部的压力为\xzanswer{D}

\fourchoices[answer=D]
{$\frac{1}{3}mg$}
{$2mg$}
{$mg$}
{$\frac{4}{3}mg$}

%% answer: D
%% memo:
\memoanswer{
对人进行受力分析，竖直向上的加速度为 $\frac{1}{3}g$，由牛顿第二定律得 $N-mg=\frac{1}{3}mg$，解得 $N=\frac{4}{3}mg$，由于人对电梯的压力与电梯对人的支持力为作用力和反作用力，故 D 正确。
}

\item
%% number: 15
%% paperName: 2005年全国理综I第 15 题 6 分
%% typeId: 2
%% type: 多选题
%% chapter: 原子和原子核
%% point: 粒子物理
%% method: 无
%% score: 6
%% degree: 650
%% duplicateId: 0
%% body:
已知 $\pi^{+}$ 介子、$\pi^{-}$ 介子都是由一个夸克（夸克 u 或夸克 d）和一个反夸克（反夸克 $\overline{\mathrm{u}}$ 或反夸克 $\overline{\mathrm{d}}$）组成的，它们的带电量如下表所示，表中 $e$ 为元电荷。

\begin{center}
\begin{tabular}{|c|c|c|c|c|c|c|}
\hline
 & $\pi^{+}$ & $\pi^{-}$ & u & d & $\overline{\mathrm{u}}$ & $\overline{\mathrm{d}}$ \\
\hline
带电量 & $+e$ & $-e$ & $+\frac{2}{3}e$ & $-\frac{1}{3}e$ & $-\frac{2}{3}e$ & $+\frac{1}{3}e$ \\
\hline
\end{tabular}
\end{center}

下列说法正确的是\xzanswer{AD}

\fourchoices[answer=AD]
{$\pi^{+}$ 由 u 和 $\overline{\mathrm{d}}$ 组成}
{$\pi^{+}$ 由 d 和 $\overline{\mathrm{u}}$ 组成}
{$\pi^{-}$ 由 u 和 $\overline{\mathrm{d}}$ 组成}
{$\pi^{-}$ 由 d 和 $\overline{\mathrm{u}}$ 组成}

%% answer: AD
%% memo:
\memoanswer{
由于 $\pi^{+}$ 介子和 $\pi^{-}$ 介子都带有一个单位的电荷电量，即 $\pi^{+}$ 介子带 $+e$、$\pi^{-}$ 介子带 $-e$ 电量，由电荷数守恒可知 A、D 正确。
}

\item
%% number: 16
%% paperName: 2005年全国理综I第 16 题 6 分
%% typeId: 2
%% type: 多选题
%% chapter: 曲线运动
%% point: 人造地球卫星
%% method: 无
%% score: 6
%% degree: 650
%% duplicateId: 0
%% body:
把火星和地球绕太阳运行的轨道视为圆周。由火星和地球绕太阳运动的周期之比可求得\xzanswer{CD}

\fourchoices[answer=CD]
{火星和地球的质量之比}
{火星和太阳的质量之比}
{火星和地球到太阳的距离之比}
{火星和地球绕太阳运行速度大小之比}

%% answer: CD
%% memo:
\memoanswer{
当行星绕太阳运动时，有 $G\frac{Mm}{R^{2}}=mR\left(\frac{2\pi}{T}\right)^{2}$，可得 $\frac{T^{2}}{R^{3}}=\frac{4\pi^{2}}{GM}=$ 恒量，故由火星与地球绕太阳的周期比可求两者环绕运行的半径比，故 C 正确；另由 $T=\frac{2\pi R}{v}$ 可得，已知两者的周期比和已求出的半径比，两者绕太阳运行的线速度比也可求解，故 CD 正确。
}

\item
%% number: 17
%% paperName: 2005年全国理综I第 17 题 6 分
%% typeId: 1
%% type: 单选题
%% chapter: 光学
%% point: 光的折射
%% method: 无
%% score: 6
%% degree: 650
%% duplicateId: 0
%% body:
图示为一直角棱镜的横截面，$\angle bac=90^\circ$，$\angle abc=60^\circ$。一平行细光束从 O 点沿垂直于 bc 面的方向射入棱镜。已知棱镜材料的折射率 $n=\sqrt{2}$，若不考虑原入射光在 bc 面上的反射光，则有光线\xzanswer{BD}

\onepicture[w=5cm]{figs/17.png}

\fourchoices[answer=BD]
{从 ab 面射出}
{从 ac 面射出}
{从 bc 面射出，且与 bc 面斜交}
{从 bc 面射出，且与 bc 面垂直}

%% answer: BD
%% memo:
\memoanswer{
光路图如图所示，由题意可求介质的临界角为 $C=\arcsin\frac{1}{n}=\arcsin\frac{1}{\sqrt{2}}=45^\circ$，故 BD 正确。

\onepicture[w=6cm]{figs/17_an.png}
}

\item
%% number: 18
%% paperName: 2005年全国理综I第 18 题 6 分
%% typeId: 1
%% type: 单选题
%% chapter: 机械波
%% point: 波长频率波速
%% method: 无
%% score: 6
%% degree: 650
%% duplicateId: 0
%% body:
一列沿 $x$ 轴正方向传播的简谐横波，周期为 $0.50\Us$。某一时刻，离开平衡位置的位移都相等的各质元依次为 P$_{1}$，P$_{2}$，P$_{3}$，$\cdots$。已知 P$_{1}$ 和 P$_{2}$ 之间的距离为 $20\Ucm$，P$_{2}$ 和 P$_{3}$ 之间的距离为 $80\Ucm$，则 P$_{1}$ 的振动传到 P$_{2}$ 所需的时间为\xzanswer{C}

\fourchoices[answer=C]
{$0.50\Us$}
{$0.13\Us$}
{$0.10\Us$}
{$0.20\Us$}

%% answer: C
%% memo:
\memoanswer{
由题意可知，P$_{1}$ 到 P$_{3}$ 之间的距离为 1 个波长，即 $\lambda=20\Ucm+80\Ucm=100\Ucm$，故 $v=\frac{\lambda}{T}=200\Ucms$，所以 P$_{1}$ 的振动传到 P$_{2}$ 时所需时间 $t=\frac{s}{v}=\frac{20\Ucm}{200\Ucms}=0.1\Us$，故 C 正确。

\onepicture[w=6cm]{figs/18_an.png}
}

\item
%% number: 19
%% paperName: 2005年全国理综I第 19 题 6 分
%% typeId: 1
%% type: 单选题
%% chapter: 电磁感应
%% point: 电磁感应中图象
%% method: 等效替代
%% score: 6
%% degree: 450
%% duplicateId: 0
%% body:
如图\ref{2005全国理综I19a}所示，两条平行虚线之间存在匀强磁场，虚线间的距离为 $l$，磁场方向垂直纸面向里。abcd 是位于纸面内的梯形线圈，ad 与 bc 间的距离也为 $l$。$t=0$ 时刻，bc 边与磁场区域边界重合（如图）。现令线圈以恒定的速度 $v$ 沿垂直于磁场区域边界的方向穿过磁场区域。取沿 a$\to$b$\to$c$\to$d$\to$a 的感应电流为正，则在线圈穿越磁场区域的过程中，感应电流 $I$ 随时间 $t$ 变化的图线可能是\xzanswer{B}

\onepicture[labelprefix={}, num={（a）}, label=2005全国理综I19a, h=4.2cm]{figs/19a.png}

\fourchoices[answer=B, ispicture=true, w=2.6cm]
{figs/19b.png}
{figs/19c.png}
{figs/19d.png}
{figs/19e.png}

%% answer: B
%% memo:
\memoanswer{
由楞次定律可得，线圈进磁场时电流方向为逆时针方向，出磁场时电流方向为顺时针方向，而无论线圈进磁场还是出磁场，线圈有效切割长度均增大，故感应电动势增大，而线圈电阻不变，感应电流增大，故 B 正确。
}

\item
%% number: 20
%% paperName: 2005年全国理综I第 20 题 6 分
%% typeId: 1
%% type: 单选题
%% chapter: 磁场
%% point: 带电粒子在磁场中的运动
%% method: 无
%% score: 6
%% degree: 450
%% duplicateId: 0
%% body:
如图，在一水平放置的平板 MN 的上方有匀强磁场，磁感应强度的大小为 $B$，磁场方向垂直于纸面向里。许多质量为 $m$ 带电量为 $+q$ 的粒子，以相同的速率 $v$ 沿位于纸面内的各个方向，由小孔 O 射入磁场区域。不计重力，不计粒子间的相互影响。下列图中阴影部分表示带电粒子可能经过的区域，其中 $R=\frac{mv}{qB}$。哪个图是正确的？\xzanswer{A}

\onepicture[w=5cm]{figs/20a.png}

\fourchoices[answer=A, ispicture=true, w=3.4cm]
{figs/20b.png}
{figs/20c.png}
{figs/20d.png}
{figs/20e.png}

%% answer: A
%% memo:
\memoanswer{
由题意可知，若粒子速度水平向右，则它在磁场中将向右做匀速圆周运动，且最远点离 O 点距离为 $2R$（竖直方向），当粒子速度方向竖直向上时，它在水平方向到达 O 点的最远距离为 $2R$，故 A 正确。
}

\item
%% number: 21
%% paperName: 2005年全国理综I第 21 题 6 分
%% typeId: 2
%% type: 多选题
%% chapter: 气体性质
%% point: 热力学第一定律
%% method: 无
%% score: 6
%% degree: 450
%% duplicateId: 0
%% body:
如图所示，绝热隔板 K 把绝热的气缸分隔成体积相等的两部分，K 与气缸壁的接触是光滑的。两部分中分别盛有相同质量、相同温度的同种气体 a 和 b。气体分子之间相互作用势能可忽略。现通过电热丝对气体 a 加热一段时间后，a、b 各自达到新的平衡\xzanswer{BCD}

\onepicture[w=5cm]{figs/21.png}

\fourchoices[answer=BCD]
{a 的体积增大了，压强变小了}
{b 的温度升高了}
{加热后 a 的分子热运动比 b 的分子热运动更激烈}
{a 增加的内能大于 b 增加的内能}

%% answer: BCD
%% memo:
\memoanswer{
由于 a 气体温度升高，压强增大，故气体将推动 K 向右运动，当达到新的平衡时，a、b 气体压强相等，此时 a 的体积增大，b 的体积减小，故 A 错误；由 $\frac{pV}{T}=$ 恒量可知，a 气体的压强增大，体积增大，故温度升高，由于 K 对气体 b 做功，故 b 气体内能增大，温度升高，故 B 正确；由于 a 的体积大于 b 的体积，且达到新的平衡后压强相同，单位体积内，a 气体的分子数小于 b 气体的分子数，所以加热后 a 的分子热运动比 b 的分子热运动更剧烈，故 C 正确；另外，由于加热前 a、b 气体温度相同，而加热后 a 的分子热运动大于 b，因此加热后 $T_{a}>T_{b}$，所以 a 增加的内能大于 b 增加的内能，故 D 正确。
}

\item
%% number: 22
%% paperName: 2005年全国理综I第 22 题 11 分
%% typeId: 4
%% type: 实验题
%% chapter: 电路
%% point: 测电动势与内阻
%% method: 无
%% score: 11
%% degree: 650
%% duplicateId: 0
%% body:
\begin{enumerate}
	\item
	在“验证力的平行四边形定则”实验中，需要将橡皮条的一端固定在水平木板上，另一端系上两根细绳，细绳的另一端都有绳套（如图）。实验中需用两个弹簧秤分别勾住绳套，并互成角度地拉橡皮条。某同学认为在此过程中必须注意以下几项：

	（A）两根细绳必须等长

	（B）橡皮条应与两绳夹角的平分线在同一直线上

	（C）在使用弹簧秤时要注意使弹簧秤与木板平面平行。

	其中正确的是 \tkanswer[2cm]{C}。

	\onepicture[w=4cm, align=right]{figs/22.png}

	\item
	测量电源 B 的电动势 $E$ 及内阻 $r$（$E$ 约为 $4.5\UV$，$r$ 约为 $1.5\UO$）。

	器材：量程 $3\UV$ 的理想电压表 V，量程 $0.5\UA$ 的电流表 A（具有一定内阻），固定电阻 $R=4\UO$，滑线变阻器 $R'$，电键 K，导线若干。

	\begin{steps}
		\item
		画出实验电路原理图。图中各元件需用题目中给出的符号或字母标出。
		\item
		实验中，当电流表读数为 $I_{1}$ 时，电压表读数为 $U_{1}$；当电流表读数为 $I_{2}$ 时，电压表读数为 $U_{2}$。则可以求出 $E=$ \tkanswer[2.5cm]{}，$r=$ \tkanswer[2.5cm]{}。（用 $I_{1}$，$I_{2}$，$U_{1}$，$U_{2}$ 及 $R$ 表示）
	\end{steps}
\end{enumerate}

%% answer:
\jdanswer{
\begin{enumerate}
	\item
	C。
	\item
	\ccnum{①}实验电路原理图如图。

	\ccnum{②}$E=\frac{I_{1}U_{2}-I_{2}U_{1}}{I_{1}-I_{2}}$，$r=\frac{U_{2}-U_{1}}{I_{1}-I_{2}}-R$。
\end{enumerate}
}

%% memo:
\memoanswer{
（1）受力与细绳的长度无关，受力只与橡皮条伸长的长度有关，故 A 错误；橡皮条与两绳夹角的角平分线在同一直线上，即合力与分力的夹角相等，这是一种特殊情况，如果合力与分力的夹角不相等，照样可以验证实验，故 B 错误；如果弹簧秤与木板平面不平行的话，即会有分力不平行木板，则其合力也不平行木板，无法证明平行四边形定则，故 C 正确。

（2）\ccnum{①}测量电源的电动势与内阻需使用限流电路，电压表为理想电压表，应将电流表内接与电压表。实验中两只电表都不能超出量程，要使电流表 A 不超出量程，即 $R_{\text{总}}\geq\frac{E}{I}=9\UO$，即外电路的总阻值必须大于 $6.5\UO$ 左右；电压表量程为 $3\UV$，而电源电动势为 $4.5\UV$，为同时保护两电表，应在电池外部接入固定电阻 $R$，设计的实验原理图如图所示。

\onepicture[w=5cm]{figs/22_an.png}

\ccnum{②}实验中，当电流表读数为 $I_{1}$ 时，电压表读数为 $U_{1}$；当电流表读数为 $I_{2}$ 时，电压表读数为 $U_{2}$，由闭合电路的欧姆定律得

$E=U_{1}+I_{1}R+I_{1}r$，$E=U_{2}+I_{2}R+I_{2}r$，

联立解得 $E=\frac{I_{1}U_{2}-I_{2}U_{1}}{I_{1}-I_{2}}$，$r=\frac{U_{2}-U_{1}}{I_{1}-I_{2}}-R$。
}

\item
%% number: 23
%% paperName: 2005年全国理综I第 23 题 16 分
%% typeId: 6
%% type: 计算题
%% chapter: 直线运动
%% point: 竖直上下抛运动
%% method: 无
%% score: 16
%% degree: 450
%% duplicateId: 0
%% body:
原地起跳时，先屈腿下蹲，然后突然蹬地。从开始蹬地到离地是加速过程（视为匀加速），加速过程中重心上升的距离称为“加速距离”。离地后重心继续上升，在此过程中重心上升的最大距离称为“竖直高度”。现有下列数据：人原地上跳的“加速距离” $d_{1}=0.50\Um$，“竖直高度” $h_{1}=1.0\Um$；跳蚤原地上跳的“加速距离” $d_{2}=0.00080\Um$，“竖直高度” $h_{2}=0.10\Um$。假想人具有与跳蚤相等的起跳加速度，而“加速距离”仍为 $0.50\Um$，则人上跳的“竖直高度”是多少？

%% answer:
\jdanswer{$H=62.5\Um$}

%% memo:
\memoanswer{
用 $a$ 表示跳蚤起跳的加速度，$v$ 表示离地时的速度，则对加速过程和离地后上升过程，分别有

$v^{2}=2ad_{2}$ \ccnum{①}，

$v^{2}=2gh_{2}$ \ccnum{②}，

若假想人具有和跳蚤相同的加速度 $a$，令 $V$ 表示在这种假想下人离地时的速度，$H$ 表示与此相应的竖直高度，则对加速过程和离地后上升过程，分别有

$V^{2}=2ad_{1}$ \ccnum{③}，

$V^{2}=2gH$ \ccnum{④}，

由以上各式可得 $H=\frac{h_{2}d_{1}}{d_{2}}$ \ccnum{⑤}，

代入数值解得 $H=62.50\Um$ \ccnum{⑥}。
}

\item
%% number: 24
%% paperName: 2005年全国理综I第 24 题 19 分
%% typeId: 6
%% type: 计算题
%% chapter: 功和能
%% point: 机械能守恒定律
%% method: 整体与隔离
%% score: 19
%% degree: 450
%% duplicateId: 0
%% body:
如图，质量为 $m_{1}$ 的物体 A 经一轻质弹簧与下方地面上的质量为 $m_{2}$ 的物体 B 相连，弹簧的劲度系数为 $k$，A、B 都处于静止状态。一条不可伸长的轻绳绕过轻滑轮，一端连物体 A，另一端连一轻挂钩。开始时各段绳都处于伸直状态，A 上方的一段绳沿竖直方向。现在挂钩上升一质量为 $m_{3}$ 的物体 C 并从静止状态释放，已知它恰好能使 B 离开地面但不继续上升。若将 C 换成另一个质量为 $(m_{1}+m_{2})$ 的物体 D，仍从上述初始位置由静止状态释放，则这次 B 刚离地时 D 的速度的大小是多少？已知重力加速度为 $g$。

\onepicture[h=5cm, align=right]{figs/24.png}

%% answer:
\jdanswer{$v=\sqrt{\frac{2m_{1}(m_{1}+m_{2})g^{2}}{(2m_{1}+m_{3})k}}$}

%% memo:
\memoanswer{
开始时，A、B 静止，设弹簧压缩量为 $x_{1}$，由胡克定律和平衡条件得 $kx_{1}=m_{1}g$ \ccnum{①}，

挂 C 并释放后，C 向下运动，A 向上运动，设 B 刚要离地时弹簧伸长量为 $x_{2}$，由受力分析知

$kx_{2}=m_{2}g$ \ccnum{②}，

B 不再上升，表示此时 A 和 C 的速度为零，C 已降到其最低点。由机械能守恒，与初始状态相比，弹簧弹性势能的增加量为

$\Delta E=m_{3}g(x_{1}+x_{2})-m_{1}g(x_{1}+x_{2})$ \ccnum{③}，

C 换成 D 后，当 B 刚离地时弹簧势能的增量与前一次相同，由能量关系得

$\frac{1}{2}(m_{3}+m_{1})v^{2}+\frac{1}{2}m_{1}v^{2}=(m_{3}+m_{1})g(x_{1}+x_{2})-m_{1}g(x_{1}+x_{2})-\Delta E$ \ccnum{④}，

由\ccnum{③}\ccnum{④}式得

$\frac{1}{2}(2m_{1}+m_{3})v^{2}=m_{1}g(x_{1}+x_{2})$ \ccnum{⑤}，

由\ccnum{①}\ccnum{②}\ccnum{⑤}式得

$v=\sqrt{\frac{2m_{1}(m_{1}+m_{2})g^{2}}{(2m_{1}+m_{3})k}}$ \ccnum{⑥}。
}

\item
%% number: 25
%% paperName: 2005年全国理综I第 25 题 20 分
%% typeId: 6
%% type: 计算题
%% chapter: 电场
%% point: 带电粒子的偏转
%% method: 无
%% score: 20
%% degree: 450
%% duplicateId: 0
%% body:
图 1 中 B 为电源，电动势 $\varepsilon=27\UV$，内阻不计。固定电阻 $R_{1}=500\UO$，$R_{2}$ 为光敏电阻。C 为平行板电容器，虚线到两极板距离相等，极板长 $l_{1}=8.0\times10^{-2}\Um$，两极板的间距 $d=1.0\times10^{-2}\Um$。S 为屏，与极板垂直，到极板的距离 $l_{2}=0.16\Um$。P 为一圆盘，由形状相同、透光率不同的三个扇形 a、b 和 c 构成，它可绕 AA$'$ 轴转动。当细光束通过扇形 a、b、c 照射光敏电阻 $R_{2}$ 时，$R_{2}$ 的阻值分别为 $1000\UO$、$2000\UO$、$4500\UO$。有一细电子束沿图中虚线以速度 $v_{0}=8.0\times10^{5}\Ums$ 连续不断地射入 C。已知电子电量 $e=1.6\times10^{-19}\UC$，电子质量 $m=9\times10^{-31}\Ukg$。忽略细光束的宽度、电容器的充电放电时间及电子所受的重力。假设照在 $R_{2}$ 上的光强发生变化时 $R_{2}$ 阻值立即有相应的改变。

\begin{enumerate}
	\item
	设圆盘不转动，细光束通过 b 照射到 $R_{2}$ 上，求电子到达屏 S 上时，它离 O 点的距离 $y$。（计算结果保留二位有效数字）。
	\item
	设转盘按图 1 中箭头方向匀速转动，每 3 秒转一圈。取光束照在 a、b 分界处时 $t=0$，试在图 2 给出的坐标纸上，画出电子到达屏 S 上时，它离 O 点的距离 $y$ 随时间 $t$ 的变化图线（$0\sim6\Us$ 间）。要求在 $y$ 轴上标出图线最高点与最低点的值。（不要求写出计算过程，只按画出的图线评分。）
\end{enumerate}

\twopicture[numA=1, numB=2, labelprefix={图},
	labelA=2005全国理综I25a, labelB=2005全国理综I25b, wA=9cm, wB=5cm]
{figs/25a.png}
{figs/25b.png}

%% answer:
\jdanswer{
\begin{enumerate}
	\item
	$y=2.4\times10^{-2}\Um$。
	\item
	图线如图所示。
\end{enumerate}
}

%% memo:
\memoanswer{
（1）设电容器 C 两板间的电压为 $U$，电场强度大小为 $E$，电子在极板间穿行时 y 方向上的加速度大小为 $a$，穿过 C 的时间为 $t_{1}$，穿出时电子偏转的距离为 $y_{1}$，由欧姆定律得

$U=\frac{\varepsilon R_{1}}{R_{1}+R_{2}}$ \ccnum{①}，

由运动学规律与牛顿第二定律得

$E=\frac{U}{d}$ \ccnum{②}，$eE=ma$ \ccnum{③}，$t_{1}=\frac{l_{1}}{v_{0}}$ \ccnum{④}，$y_{1}=\frac{1}{2}at_{1}^{2}$ \ccnum{⑤}，

联立解得 $y_{1}=\frac{e\varepsilon R_{1}l_{1}^{2}}{2dmv_{0}^{2}(R_{1}+R_{2})}$ \ccnum{⑥}，

代入数据得

$y_{1}=4.8\times10^{-3}\Um$ \ccnum{⑦}，

由此可见 $y_{1}<\frac{1}{2}d$，电子可通过 C。

设电子从 C 穿出时，沿 y 方向的速度为 $v_{y}$，穿出后到达屏 S 所经历的时间为 $t_{2}$，在此时间内电子在 y 方向移动的距离为 $y_{2}$，则有

$v_{y}=at_{1}$ \ccnum{⑧}，$t_{2}=\frac{l_{2}}{v_{0}}$ \ccnum{⑨}，$y_{2}=v_{y}t_{2}$ \ccnum{⑩}，

联立解得 $y_{2}=\frac{e\varepsilon R_{1}l_{1}l_{2}}{dmv_{0}^{2}(R_{1}+R_{2})}$ \ccnum{⑪}，

代入数据得

$y_{2}=1.92\times10^{-2}\Um$ \ccnum{⑫}，

由题意 $y=y_{1}+y_{2}=2.4\times10^{-2}\Um$ \ccnum{⑬}。

（2）如图所示（$R_{2}=1000\UO$ 时，电子未通过 C）。

\onepicture[w=5cm]{figs/25_an.png}
}

\end{enumerate}

\end{document}
